\documentclass[11pt]{article}

\usepackage[margin=1in]{geometry}
\usepackage{amsmath,amssymb}
\usepackage{bm}
\usepackage{hyperref}

\title{A Simple KFBIM Shape Optimization Problem}
\author{KFBIM notes}
\date{\today}

\begin{document}
\maketitle

\section*{Goal}

The application \texttt{shape\_opt\_transmission\_2d} solves a small inverse
shape problem for a two-phase screened diffusion equation.  The unknown is the
boundary $\Gamma(a)$ of an inclusion $D(a)\subset \Omega$, where
$\Omega=(-1.2,1.2)^2$.  Synthetic target data are generated from a known hidden
shape, and the optimization attempts to recover that shape from exterior field
measurements.

\section*{Forward Transmission Problem}

For each outer boundary excitation $g_\ell$, solve
\begin{align}
    -\nabla\cdot(\beta \nabla u_\ell) + \kappa^2 u_\ell &= 0
        && \text{in } \Omega\setminus\Gamma(a), \\
    u_\ell &= g_\ell
        && \text{on } \partial\Omega, \\
    [u_\ell] &= 0
        && \text{on } \Gamma(a), \\
    [\beta\,\partial_n u_\ell] &= 0
        && \text{on } \Gamma(a).
\end{align}
The coefficient is piecewise constant:
\[
    \beta =
    \begin{cases}
    \beta_{\rm int}, & x\in D(a),\\
    \beta_{\rm ext}, & x\in \Omega\setminus\overline{D(a)}.
    \end{cases}
\]
The implemented demo uses
\[
    \beta_{\rm int}=4,\qquad
    \beta_{\rm ext}=1,\qquad
    \frac{\kappa_{\rm int}^2}{\beta_{\rm int}}
    =
    \frac{\kappa_{\rm ext}^2}{\beta_{\rm ext}}
    =
    \lambda^2=1.1.
\]
Thus the KFBIM common-ratio transmission API can reduce both phases to the
same screened operator $-\Delta+\lambda^2$ after division by $\beta$.

The boundary excitations are
\[
    g_0=1,\qquad
    g_1=x,\qquad
    g_2=y,\qquad
    g_3=x^2-y^2.
\]
For a candidate shape, KFBIM returns the grid solution values
$u_\ell(x_m;a)$ at fixed observation nodes $x_m$ in an exterior annulus.

\section*{Shape Parameterization}

The inclusion boundary is represented as a fixed-center radial Fourier curve:
\[
    \Gamma(a):\quad
    X(\theta;a)
    =
    c + R(\theta;a)(\cos\theta,\sin\theta),
    \qquad
    c=(0.07,-0.04),
\]
with
\[
    R(\theta;a)
    =
    R_0\exp\left(
        \sum_{k=1}^{K}
        a_k\cos(k\theta)+b_k\sin(k\theta)
    \right),
    \qquad R_0=0.5,\quad K=3.
\]
The exponential form keeps the radius positive.  The hidden target used by the
demo is
\[
    (a_1,b_1,a_2,b_2,a_3,b_3)
    =
    (0.04,0,0,-0.08,0.16,0).
\]

\section*{Optimization Objective}

Let $d_{\ell m}=u_\ell(x_m;a_\star)$ denote synthetic measurements from the
target shape $a_\star$.  The demo minimizes
\[
    J(a)
    =
    \frac{1}{2M}
    \sum_{\ell,m}
    \left(u_\ell(x_m;a)-d_{\ell m}\right)^2
    + \alpha
    \sum_{k=1}^{K} k^4(a_k^2+b_k^2)
    + \rho\left(|D(a)|-|D(a_\star)|\right)^2.
\]
The first term measures mismatch to the target data.  The second term is a
smoothness regularizer that penalizes high Fourier modes.  The third term
discourages area drift.

\section*{Adjoint Shape Gradient}

The optimized application can replace finite-difference gradients by one
continuous adjoint solve per boundary excitation.  For each forward solution
$u_\ell$, the adjoint $p_\ell$ solves
\[
    -\nabla\cdot(\beta\nabla p_\ell)+\kappa^2p_\ell
    =
    \frac{1}{M_{\rm tot}}\sum_m
    \left(u_\ell(x_m;a)-d_{\ell m}\right)\delta_{x_m},
\]
with homogeneous Dirichlet data on $\partial\Omega$ and the same transmission
conditions
\[
    [p_\ell]=0,\qquad [\beta\,\partial_n p_\ell]=0
    \quad\text{on }\Gamma(a).
\]
On the Cartesian grid, each point source is represented by a nodal delta with
weight proportional to $1/h^2$.  Since the KFBIM transmission API takes the
reduced right-hand side $q=f/\beta$, the exterior observation source is entered
as
\[
    q(x_m)=
    \frac{u_\ell(x_m;a)-d_{\ell m}}
         {M_{\rm tot}\,\beta_{\rm ext}\,h^2}.
\]

Let $n$ be the normal pointing from the inclusion to the exterior and let
$\tau$ be a unit tangent.  Because both the state and adjoint are continuous
across the interface, their tangential derivatives are single-valued.  The
implemented continuous shape-gradient density is
\[
    g_\Gamma
    =
    \sum_\ell
    \left\{
    \beta_{\rm int}\,\partial_n u_{\ell,{\rm int}}\,
        \partial_n p_{\ell,{\rm int}}
    -
    \beta_{\rm ext}\,\partial_n u_{\ell,{\rm ext}}\,
        \partial_n p_{\ell,{\rm ext}}
    -
    (\beta_{\rm int}-\beta_{\rm ext})
    \left[
        \partial_\tau u_\ell\,\partial_\tau p_\ell
        + \lambda^2 u_\ell p_\ell
    \right]
    \right\}.
\]
For a Fourier coefficient, this density is projected onto the normal velocity
generated by that coefficient.  For example,
\[
    \frac{\partial J_{\rm data}}{\partial a_k}
    =
    \int_\Gamma
    g_\Gamma\,
    R(\theta)\cos(k\theta)\,
    (e_r\cdot n)\,ds,
\]
and the formula for $b_k$ replaces $\cos(k\theta)$ by $\sin(k\theta)$.  The
regularization and area-penalty derivatives are then added explicitly.

\section*{Numerical Method}

For each trial vector $a$, the code:
\begin{enumerate}
    \item builds a P2 quadratic interface with
          \texttt{CurveResampler2D::discretize};
    \item solves the four common-ratio transmission problems using
          \texttt{LaplaceTransmission2D};
    \item samples the resulting fields at the observation nodes;
    \item evaluates $J(a)$.
\end{enumerate}
The default implementation now uses the adjoint gradient with a backtracking
gradient-descent step.  The finite-difference gradient remains available for
verification.  A gradient check compares the adjoint direction against central
finite differences at the current shape.

\section*{Important Practical Point}

If the regularization weight $\alpha$ is too large, the exact target shape may
not minimize the chosen objective.  In the default target, the $k=3$ mode is
important, but the regularizer weights it by $3^4=81$.  For clean synthetic
recovery, use a very small $\alpha$ or set $\alpha=0$.  For noisy data, keep a
small positive $\alpha$ to suppress nonphysical high-frequency shape changes.

\section*{Research-Level Direction}

The demo is intentionally small, but it points to a larger class of inverse
interface problems that is well matched to KFBIM: recover a moving embedded
biophysical or materials boundary from fields measured away from the boundary.
One concrete story is an active cell or organelle whose membrane separates two
diffusive media with different transport coefficients.  Boundary actuation or
external illumination produces several steady screened-diffusion fields, while
microscopy or electrodes measure only the exterior response.  The inverse
problem is to infer the hidden membrane shape, not merely the scalar
coefficient values.

This setting is difficult for a standard body-fitted finite-element method
because every shape update changes the mesh and the interface geometry.  It is
also awkward for a classical boundary integral equation when the bulk operator,
forcing, boundary conditions, or future variable coefficients destroy the
availability of a simple analytic Green's function.  KFBIM keeps the interface
geometric unknown explicit while using Cartesian bulk solvers, so shape
iterations can move the boundary without remeshing the full domain.

\section*{KFBIM-Specific Metric}

The current app uses a Euclidean gradient in the six Fourier coefficients.
A more KFBIM-specific iteration would replace this by a geometry-aware
preconditioned gradient.  Let $V_i(\theta)$ be the normal velocity generated by
parameter $a_i$ or $b_i$.  Instead of stepping with
\[
    \delta a_i = -\eta\,\frac{\partial J}{\partial a_i},
\]
define a Gram matrix
\[
    G_{ij}
    =
    \int_\Gamma V_i V_j\,ds
    +
    \ell^2\int_\Gamma \partial_s V_i\,\partial_s V_j\,ds
    +
    \gamma\,
    \langle \mathcal{K}V_i,\mathcal{K}V_j\rangle_{\Omega_h}.
\]
Here $\mathcal{K}V$ denotes the KFBIM linearized field response produced by a
normal boundary displacement $V$.  In practice, this response can be estimated
with the same spread--bulk-solve--restrict machinery that already applies the
PDE operator.  The update is then
\[
    G\,\delta a = -\nabla_a J.
\]

The first two terms give the usual $L^2$ or Sobolev geometry metric on the
boundary.  The last term is problem-dependent: two boundary motions are close
if KFBIM predicts that they produce similar observable PDE fields.  This can
reduce wasted steps in nearly invisible shape directions, balance Fourier
modes according to the actual measurement map rather than an arbitrary
coefficient scale, and make line search less sensitive to the grid resolution
or the chosen parameterization.

\section*{Why This Can Be Justified}

The metric above is a Gauss--Newton or natural-gradient approximation for the
reduced objective.  If $\mathcal{F}(a)$ maps shape parameters to all observed
data, then the data term is
\[
    J_{\rm data}(a)
    =
    \frac12\|\mathcal{F}(a)-d\|^2.
\]
Near a local minimizer,
\[
    \mathcal{F}(a+\delta a)
    \approx
    \mathcal{F}(a)+\mathcal{F}'(a)\delta a,
\]
so the local quadratic model has Hessian
\[
    \mathcal{F}'(a)^*\mathcal{F}'(a)
\]
plus regularization terms, up to residual-weighted second derivatives.  The
KFBIM response term in $G$ is a discretized version of this
pullback metric.  Under the usual differentiability assumptions for smooth
elliptic transmission shape functionals, this produces a descent direction
whenever $G$ is symmetric positive definite and the adjoint gradient is
consistent with the chosen discretization.  It is not a magic convexification,
but it is a principled preconditioner for the reduced-space optimization.

\end{document}
